%! Author = sriadityavedantam
%! Date = 3/25/26

\lesson{27}{Monday 13 October 2025 9:10}{Day 27}

\begin{theorem}{
    Any non-empty perfect set is uncountable.
}
    Suppose $P\neq\varnothing$ is perfect and let $x_1,x_2,\ldots$ be any sequence of points in $P$.
    We will show
    \[
        P\setminus\{x_1,x_2,\ldots\}\neq\varnothing.
    \]
    We construct closed intervals
    \[
        J_1 \supset J_2 \supset J_3 \supset \cdots
    \]
    such that:
    \begin{enumerate}
        \item $J_n\cap P\neq\varnothing$,
        \item $x_n\notin J_n$,
        \item $J_{n+1}\subset J_n$.
    \end{enumerate}

    Since $P$ is perfect, every point of $P$ is a limit point of $P$.
    Choose $z_1\in P$ with $z_1\neq x_1$.
    Let
    \[
        \eps_1 \coloneqq \frac{1}{2}\abs{z_1-x_1}.
    \]
    Set
    \[
        J_1 \coloneqq [z_1-\eps_1,z_1+\eps_1].
    \]
    Then $x_1\notin J_1$.
    Also $J_1\cap P\neq\varnothing$ since $z_1\in J_1\cap P$.
    Suppose $J_n$ has been chosen with $J_n\cap P\neq\varnothing$ and $x_n\notin J_n$.
    Since $J_n\cap P$ is non-empty and every point of $P$ is a limit point of $P$, we may choose
    \[
        z_{n+1}\in (J_n\cap P)\setminus\{x_{n+1}\}.
    \]
    Because $J_n$ is a closed interval and $z_{n+1}\in J_n$, there exists $\delta_{n+1}>0$ such that
    \[
        [z_{n+1}-\delta_{n+1},z_{n+1}+\delta_{n+1}] \subset J_n.
    \]
    Let
    \[
        \eps_{n+1} \coloneqq \frac{1}{2}\min\{\abs{z_{n+1}-x_{n+1}},\delta_{n+1}\}.
    \]
    Define
    \[
        J_{n+1} \coloneqq [z_{n+1}-\eps_{n+1},z_{n+1}+\eps_{n+1}].
    \]
    Then $J_{n+1}\subset J_n$ and $x_{n+1}\notin J_{n+1}$.
    Also $J_{n+1}\cap P\neq\varnothing$ since $z_{n+1}\in J_{n+1}\cap P$.
    Thus the construction is complete.
    Since each $J_n$ is closed and bounded, it is compact.
    Since $P$ is closed, each $J_n\cap P$ is compact.
    Also
    \[
        J_1\cap P \supset J_2\cap P \supset J_3\cap P \supset \cdots
    \]
    is a nested sequence of non-empty compact sets.
    Hence
    \[
        \bigcap_{n=1}^\infty (J_n\cap P)\neq\varnothing.
    \]
    Let
    \[
        y\in \bigcap_{n=1}^\infty (J_n\cap P).
    \]
    Then $y\in P$.
    Also, for each $n$, since $x_n\notin J_n$ and $y\in J_n$, we have $y\neq x_n$.
    Therefore
    \[
        y\in P\setminus\{x_1,x_2,\ldots\}.
    \]
    Hence no sequence can exhaust $P$.
    Therefore $P$ is uncountable.
\end{theorem}

\begin{definition}[Limit of a Function]
    Let $A\subset\r$, $x_0$ be a limit point of $A$, and $f:A\to\r$.
    We say
    \[
        \lim_{x\to x_0} f(x) = L
    \]
    if for every $\eps > 0$, there exists $\delta > 0$ such that
    \[
        0 < \abs{x-x_0} < \delta,\ x\in A
    \]
    implies
    \[
        \abs{f(x)-L} < \eps.
    \]
\end{definition}